\documentclass[11pt]{amsart}
\usepackage{amsmath,amssymb,amsthm,mathtools}
\begin{document}

Throughout, all digit and carry conditions are in base $p$.
For nonnegative integers $x_1,\ldots,x_r$, we write
$x_1\oplus x_2\oplus\cdots\oplus x_r$ to indicate that the addition
$x_1+x_2+\cdots+x_r$ has no carries in base $p$; equivalently,
at every base-$p$ digit position, the digits of $x_1,\ldots,x_r$ sum to at most $p-1$.

For integers $d\geq 1$ and $j\geq 0$, let $\mathcal{T}_{d,j}$ denote the set of
$d$-tuples $(m_1,\ldots,m_d)$ of positive multiples of $q-1$ such that
$j\oplus m_1\oplus\cdots\oplus m_d$ holds, and define
\[
  M_d(j) \coloneq \min_{(m_1,\ldots,m_d)\in\mathcal{T}_{d,j}}\left(m_1+2m_2+\cdots+dm_d\right),
\]
so that $s_d(k)=dk+M_d(k-1)$.

By definition,
\[
  s_d(k)=dk+M_d(k-1),
  \qquad
  s_d(k+1)=d(k+1)+M_d(k).
\]
Therefore
\begin{equation}\label{eq:H3-difference}
  s_d(k+1)-s_d(k)=d+M_d(k)-M_d(k-1).
\end{equation}
If $p\nmid k$, then the units digit of $k$ in base-$p$ is nonzero.
Passing from $k$ to $k-1$ lowers only that units digit by $1$ and leaves all higher digits unchanged.
Hence
\[
  k\oplus m_1\oplus\cdots\oplus m_d \text{ carry-free}
  \implies (k-1)\oplus m_1\oplus\cdots\oplus m_d \text{ carry-free}.
\]
This implies $\mathcal T_{d,k} \subseteq \mathcal T_{d,k-1}$ and we have
\[ M_d(k)\geq M_d(k-1).  \]
Hence \eqref{eq:H3-difference} gives $s_d(k+1)-s_d(k)\geq d>0$ and we're done.

\end{document}
